% !PRINTED none
% Abschnitt title page: no running head/page number printed on this page
% (chapter-opening convention); full body text transcribed below.

\begin{center}
\gesperrt{Neunzehnter Abschnitt.}

\vspace{1em}
{\Large Berechnung der Klasseninvarianten.}

\vspace{0.5em}
\rule{2cm}{0.4pt}

\vspace{1em}
\S\,125. \gesperrt{Die Klasseninvariante $\gamma_2$.}
\end{center}

Wir haben schon oben bemerkt, daß außer der Funktion $j(\omega)$ noch
andere Modulfunktionen zur Bildung von Klasseninvarianten geeignet
sind, und oft zu einfacheren Resultaten Anlaß geben. Unter diesen
tritt uns zunächst die Funktion $\gamma_2(\omega)$ entgegen, die durch
$\sqrt[3]{j(\omega)}$ definirt ist. % [sic]

Wir haben im \S\,71 gesehen, daß zwischen den Funktionen

\weq{1}{u = \gamma_2(\omega), \quad v = \gamma_2\left(\frac{c+\partial\omega}{a}\right),}

falls $a\partial = n$ durch 3 nicht teilbar ist, und $c$ durch 3
teilbar angenommen ist, eine Transformationsgleichung

\weq{2}{\Phi_n(u, v) = 0}

besteht. Die Funktion $\Phi_n$ hängt aber nur von den beiden
Argumenten $vu^{-n}$, $u^3$ ab, und in $\Phi_n(u, u) = 0$ kommt also,
wenn $n \equiv 1\ (\mathrm{mod}\ 3)$ ist, nur $u^3$, d.\,h. $j(\omega)$
vor; in diesem Falle kann daher diese Gleichung kein anderes Resultat
ergeben, als die Invariantengleichung. Anders ist es aber in dem
Falle

\weq{3}{n \equiv -1\ (\mathrm{mod}\ 3),}

den wir jetzt voraussetzen wollen.

Wenn einer der Werte $v = u$ werden soll, so muß eine Relation
bestehen

\weq{4}{\frac{\gamma+\delta\omega}{\alpha+\beta\omega} = \frac{c+\partial\omega}{a}, \quad \alpha\delta-\beta\gamma = 1,}

und zugleich muß [\S\,54, (15)]

\weq{5}{-\alpha\beta+\alpha\gamma+\beta\delta-\alpha^2\beta\delta \equiv 0\ (\mathrm{mod}\ 3)}

sein. Es ist also jedenfalls $\omega$ Wurzel einer quadratischen
Gleichung:

\weq{6}{A\omega^2+B\omega+C = 0, \quad D = B^2-4AC.}
